\subsection*{6.1 Определение и вопрос существования}

В шестой главе мы познакомимся с таким непростым объектом теории вероятностей, как условное математическое ожидание, а также научимся некоторым способам его вычисления.

Мы хорошо помним понятие условной вероятности из главы 1. Пусть теперь $\xi$ и $\eta$ --- это две случайные величины. Как можно было бы определить $\E(\xi|\eta = y)$ и $\E(\xi|\eta)$? Посмотрим сначала, как можно было бы подойти к этому чисто интуитивно.

Рассмотрим сначала самый понятный случай, когда $\xi$ и $\eta$ являются простыми случайными величинами. Как мы помним, для них обычное математическое ожидание определяется по формуле
$$ \E\xi = \sum_{i=1}^n x_i\P(\xi = x_i), $$
где $x_1, \dots, x_n$ --- это все возможные значения случайной величины $\xi$. Тогда очень естественно определить $\E(\xi|\eta = y)$ следующим образом:
\begin{equation} \label{eq:cme_discrete}
    \E(\xi|\eta = y) = \sum_{i=1}^n x_i\P(\xi = x_i|\eta = y),
\end{equation}
где $\P(\xi = x_i|\eta = y)$ --- это стандартная условная вероятность, введенная еще в первой главе.

Как можно обобщить определение \eqref{eq:cme_discrete} на случай непростой с.в. $\xi$? Домножим правую часть равенства \eqref{eq:cme_discrete} на $\P(\eta = y)$:
$$ \sum_{i=1}^n x_i\P(\xi = x_i, \eta = y) = \E\left( \sum_{i=1}^n x_i \cdot \I\{\xi = x_i\}\I\{\eta = y\} \right) = \E(\xi \cdot \I\{\eta = y\}). $$

С другой стороны, если обозначить $\E(\xi|\eta = y) = \varphi(y)$, то получим равенство
$$ \varphi(y)\P(\eta = y) = \E(\varphi(y)\I\{\eta = y\}) = \E(\varphi(\eta)\I\{\eta = y\}). $$
Данное тождество осмысленно и для непростой случайной величины $\xi$. Тем самым, разумно сказать, что $\E(\xi|\eta = y)$ --- это такая функция $\varphi(y)$, что для любого $y$ выполнено равенство
$$ \E(\xi \cdot \I\{\eta = y\}) = \E(\varphi(\eta)\I\{\eta = y\}). $$

Осталось понять, как можно было бы теперь перенести данное определение на случай уже непростой с.в. $\eta$. Для этого мы изучим более общее понятие условного математического ожидания случайной величины относительно сигма-алгебры.

Пусть $(\Omega, \mathcal{F}, \P)$ --- вероятностное пространство, $\xi$ --- случайная величина на нем, а $\mathcal{C} \subset \mathcal{F}$ --- под-$\sigma$-алгебра в $\mathcal{F}$.

\mkdef{Условное математическое ожидание относительно $\sigma$-алгебры}{cme_sigma}{
    Условным математическим ожиданием случайной величины $\xi$ относительно $\sigma$-алгебры $\mathcal{C}$ называется случайная величина, обозначаемая $\E(\xi|\mathcal{C})$, которая удовлетворяет следующим двум свойствам:
    \begin{enumerate}
        \item $\E(\xi|\mathcal{C})$ является $\mathcal{C}$-измеримой (свойство измеримости);
        \item для любого $A \in \mathcal{C}$ выполняется равенство $\E(\xi \cdot \I_A) = \E(\E(\xi|\mathcal{C})\I_A)$ или
        $$ \int\limits_A \xi \, d\P = \int\limits_A \E(\xi|\mathcal{C}) \, d\P $$
        (интегральное свойство).
    \end{enumerate}
}

Ясно, что в общей ситуации случайная величина $\xi$ не является $\mathcal{C}$-измеримой, потому $\E(\xi|\mathcal{C})$, вообще говоря, является нетривиальным объектом. Но глядя на определение сразу возникает естественный вопрос о существовании условного математического ожидания. Для ответа на него нам понадобятся некоторые факты из теории меры.

\mkdef{Конечный заряд}{charge}{
    Пусть $(\Omega, \mathcal{F}, \P)$ --- вероятностное пространство. Функция $\nu \colon \mathcal{F} \to \mathbb{R}$ называется (конечным) зарядом (или мерой со знаком), если $\nu$ --- счетно-аддитивна на $\mathcal{F}$, т.е.
    $$ \nu\left(\bigcup_{n=1}^\infty A_n\right) = \sum_{n=1}^\infty \nu(A_n), $$
    где ряд в правой части сходится абсолютно, и
    $$ \sup_{A \in \mathcal{F}} |\nu(A)| < +\infty. $$
    Заряд $\nu$ является абсолютно непрерывным относительно $\P$, если из того, что $\P(A) = 0$ следует, что $\nu(A) = 0$.
}

Следующая теорема Радона--Никодима дает хорошее представление о том, как выглядят заряды, абсолютно непрерывные относительно меры $\P$.

\mkthm{Радон, Никодим}{radon_nikodym}{
    Пусть $(\Omega, \mathcal{F}, \P)$ --- вероятностное пространство, а $\nu$ --- заряд, абсолютно непрерывный относительно $\P$. Тогда существует единственная (с точностью до равенства п.н.) такая случайная величина $\eta \in L^1(\Omega, \mathcal{F}, \P)$, что для любого $A \in \mathcal{F}$ выполнено
    $$ \nu(A) = \int\limits_A \eta \, d\P = \E(\eta \cdot \I_A). $$
}

\mkrem{
    В этом случае $\eta = \frac{d\nu}{d\P}$ называется производной Радона--Никодима.
}

Доказательство теоремы \ref{thm:radon_nikodym} можно найти, например, в [2]. Из нее легко вытекает существование условного математического ожидания в предположении конечности обычного.

\mklem{О существовании УМО}{cme_exist}{
    Если $\E|\xi| < \infty$, то для любой $\sigma$-алгебры $\mathcal{C} \subset \mathcal{F}$ условное математическое ожидание $\E(\xi|\mathcal{C})$ существует и единственно с точностью до равенства почти наверное.
}
\mkproof{
    Рассмотрим вероятностное пространство $(\Omega, \mathcal{C}, \P)$. Для любого $A \in \mathcal{C}$ положим
    $$ \nu(A) = \int\limits_A \xi \, d\P = \E(\xi \cdot \I_A). $$
    Это заряд на $(\Omega, \mathcal{C}, \P)$, абсолютно непрерывный относительно $\P$. По теореме Радона--Никодима существует единственная (п.н.) случайная величина $\eta \in L^1(\Omega, \mathcal{C}, \P)$ т.,ч. для любого $A \in \mathcal{C}$
    $$ \nu(A) = \int\limits_A \eta \, d\P = \E(\eta \cdot \I_A). $$
    Случайная величина $\eta$ является $\mathcal{C}$-измеримой и удовлетворяет интегральному свойству. Значит, по определению $\eta = \E(\xi|\mathcal{C})$.
}

\mkrem{
    Условное математическое ожидание определяется с точностью до равенства п.н. Иногда бывает важным осуществлять правильный подбор его ``варианта''. Всюду далее мы будем опускать равенство п.н., предполагая, однако, что оно выполнено, если в одной из его частей стоит УМО.
}

Лемма \ref{lem:cme_exist} отвечает на вопрос о существовании условного математического ожидания, но не дает представления о том, чему оно может равняться. Но несложно увидеть, что в случае, когда $\sigma$-алгебра $\mathcal{C}$ порождена некоторым не более чем счетным разбиением $\{D_n, n \in \mathbb{N}\}$ пространства $\Omega$, можно дать явный вид $\E(\xi|\mathcal{C})$, который в силу единственности можно даже считать эквивалентным определением.

\mklem{Явный вид для разбиения}{cme_partition}{
    Если $\sigma$-алгебра $\mathcal{C}$ порождена разбиением $\{D_n, n \in \mathbb{N}\}$, то для любой случайной величины $\xi$ с условием $\E|\xi| < +\infty$ выполнено
    $$ \E(\xi|\mathcal{C})(\omega) = \sum_{n=1}^\infty \frac{\E(\xi \cdot \I_{D_n})}{\P(D_n)} \cdot \I_{D_n}(\omega). $$
}
\mkproof{
    Отметим, что если $\omega \in D_n$ и $\P(D_n) = 0$, то правая часть формально не определена. Но это происходит с вероятностью нуль, поэтому можно считать, что для таких $\omega$ величина $\E(\xi|\mathcal{C})(\omega)$ также равна нулю.

    Наш план доказательства равенства прост. Нам в правой части предлагают кандидата в УМО, в силу единственности надо проверить для него два свойства: свойство измеримости и интегральное свойство. Правая часть есть линейная комбинация несовместных индикаторов событий $D_n \in \mathcal{C}$, так что это $\mathcal{C}$-измеримая случайная величина.

    Проверим интегральное свойство. Если $A \in \mathcal{C}$, то оно представляет собой лишь дизъюнктное объединение элементов $D_n$. В связи с этим достаточно рассмотреть ситуацию $A = D_m$ для некоторого $m \in \mathbb{N}$. Если $\P(D_m) = 0$, то искомое равенство очевидно. В остальных случаях имеем:
    \begin{align*}
        \E\left( \left( \sum_{n=1}^\infty \frac{\E(\xi \cdot \I_{D_n})}{\P(D_n)} \cdot \I_{D_n} \right) \I_{D_m} \right) &= \\
        \intertext{(пользуемся несовместностью событий $D_n$)}
        = \E\left( \frac{\E(\xi \cdot \I_{D_m})}{\P(D_m)} \cdot \I_{D_m} \right) = \frac{\E(\xi \cdot \I_{D_m})}{\P(D_m)} \cdot \P(D_m) &= \E(\xi \cdot \I_{D_m}).
    \end{align*}
}

Формула неформально показывает, что условное математическое ожидание $\E(\xi|\mathcal{C})$ является некоторым ``усреднением'' значений $\xi$ по множествам из $\mathcal{C}$.

Перейдем к обсуждению свойств условного математического ожидания.


\subsection*{6.2 Основные свойства УМО}

Всюду в теореме ниже предполагаем конечность математических ожиданий у рассматриваемых случайных величин, если это необходимо.

\mkthm{Основные свойства УМО}{cme_props}{
    \begin{enumerate}[label=(\arabic*), leftmargin=*]
        \item Если с.в. $\xi$ является $\mathcal{C}$-измеримой, то $\E(\xi|\mathcal{C}) = \xi$.
        \item Линейность. Для любых случайных величин $\xi$, $\eta$, $a, b \in \mathbb{R}$ и $\sigma$-алгебры $\mathcal{C}$ выполнено
        $$ \E(a\xi + b\eta|\mathcal{C}) = a\E(\xi|\mathcal{C}) + b\E(\eta|\mathcal{C}). $$
        \item Формула полной вероятности.
        $$ \E(\E(\xi|\mathcal{C})) = \E\xi. $$
        \item Если $\xi$ независима с $\mathcal{C}$, то $\E(\xi|\mathcal{C}) = \E\xi$.
        \item Сохранение отношения порядка. Если $\xi \ge \eta$ п.н., то
        $$ \E(\xi|\mathcal{C}) \ge \E(\eta|\mathcal{C}) \text{ п.н.} $$
        \item $$ |\E(\xi|\mathcal{C})| \le \E(|\xi| \mid \mathcal{C}). $$
        \item Телескопическое. Пусть $\xi$ --- с.в., а $\mathcal{C}_1 \subset \mathcal{C}_2$ --- две $\sigma$-алгебры. Тогда
        \begin{enumerate}[label=(\alph*)]
            \item $\E(\E(\xi|\mathcal{C}_1)|\mathcal{C}_2) = \E(\xi|\mathcal{C}_1)$,
            \item $\E(\E(\xi|\mathcal{C}_2)|\mathcal{C}_1) = \E(\xi|\mathcal{C}_1)$.
        \end{enumerate}
        \item Предельный переход под знаком УМО. Пусть $\{\xi_n, n \in \mathbb{N}\}$, $\xi$ --- случайные величины.
        \begin{enumerate}[label=(\alph*)]
            \item Если $0 \le \xi_n \uparrow \xi$ п.н., то $\E(\xi_n|\mathcal{C}) \uparrow \E(\xi|\mathcal{C})$ п.н.
            \item Пусть $\xi_n \toAS \xi$ и $|\xi_n| \le \delta$ для всех $n$, причем $\E\delta < +\infty$. Тогда 
            $$ \E(\xi_n|\mathcal{C}) \toAS \E(\xi|\mathcal{C}). $$
        \end{enumerate}
        \item Пусть $\eta$ --- это $\mathcal{C}$-измеримая с.в., то
        $$ \E(\xi \cdot \eta|\mathcal{C}) = \eta\E(\xi|\mathcal{C}). $$
        \item Неравенство Йенсена. Пусть $\varphi$ --- выпуклая книзу функция. Тогда
        $$ \E(\varphi(\xi)|\mathcal{C}) \ge \varphi(\E(\xi|\mathcal{C})). $$
    \end{enumerate}
}

\mkproof{
    \begin{enumerate}[label=\arabic*., leftmargin=*]
        \item Очевидно, что для $\xi$ выполняются и свойство измеримости, и интегральное свойство.
        \item Для начала заметим, что и $\E(\xi|\mathcal{C})$, и $\E(\eta|\mathcal{C})$ являются $\mathcal{C}$-измеримыми случайными величинами. Тогда $a\E(\xi|\mathcal{C}) + b\E(\eta|\mathcal{C})$ --- это тоже $\mathcal{C}$-измеримая случайная величина. Осталось проверить интегральное свойство. В силу линейности математического ожидания для любого $A \in \mathcal{C}$
        $$ \E((a\xi + b\eta) \cdot \I_A) = a\E(\xi \cdot \I_A) + b\E(\eta \cdot \I_A). $$
        Согласно интегральному свойству:
        $$ a\E(\xi \cdot \I_A) + b\E(\eta \cdot \I_A) = a\E(\E(\xi|\mathcal{C})\I_A) + b\E(\E(\eta|\mathcal{C})\I_A). $$
        Свернём сумму назад:
        $$ a\E(\E(\xi|\mathcal{C})\I_A) + b\E(\E(\eta|\mathcal{C})\I_A) = \E\left( (a\E(\xi|\mathcal{C}) + b\E(\eta|\mathcal{C})) \cdot \I_A \right). $$
        Тем самым, получаем желаемое.
        \item Воспользуемся интегральным свойством, положив $A = \Omega$:
        $$ \E\xi = \E(\xi \cdot \I_A) = \E(\E(\xi|\mathcal{C}) \cdot \I_A) = \E(\E(\xi|\mathcal{C})). $$
        \item Напомним, что $\xi$ независима с $\mathcal{C}$ тогда и только тогда, когда для любого $A \in \mathcal{C}$ с.в. $\xi$ независима с $\I_A$. Заметим, что $\E\xi$ --- это константа, поэтому $\E\xi$ является $\mathcal{C}$-измеримой. Проверим интегральное свойство: для любого $A \in \mathcal{C}$
        $$ \E(\xi \cdot \I_A) = [\text{независимость}] = \E\xi \cdot \E(\I_A) = \E(\E\xi \cdot \I_A). $$
        \item Отметим снова, что $\E(\xi|\mathcal{C})$ и $\E(\eta|\mathcal{C})$ являются случайными величинами на вероятностном пространстве $(\Omega, \mathcal{C}, \P)$. Для любого $A \in \mathcal{C}$ имеем:
        $$ \E(\E(\xi|\mathcal{C})\I_A) = \E(\xi \cdot \I_A) \le \E(\eta \cdot \I_A) = \E(\E(\eta|\mathcal{C})\I_A). $$
        Согласно свойству 3.9 обычных математических ожиданий получаем, что с вероятностью 1 выполнено искомое неравенство $\E(\xi|\mathcal{C}) \le \E(\eta|\mathcal{C})$.
        \item Сразу следует из свойств (2) и (5).
        \item Первый пункт почти очевиден. Заметим, что с.в. $\E(\xi|\mathcal{C}_1)$ является $\mathcal{C}_1$-измеримой, а, значит, и $\mathcal{C}_2$-измеримой. Далее, применяем свойство (1). 
        
        Теперь докажем второе утверждение. Заметим, что свойство измеримости у нас уже есть. Проверим интегральное свойство. Для любого $A \in \mathcal{C}_1 \subset \mathcal{C}_2$:
        \begin{align*}
            \E(\E(\xi|\mathcal{C}_1)\I_A) &= [\text{инт. свойство для } \mathcal{C}_1] = \E(\xi \cdot \I_A) = \\
            &= [\text{инт. свойство для } \mathcal{C}_2] = \E(\E(\xi|\mathcal{C}_2)\I_A).
        \end{align*}
        \item без доказательства
        \item Заметим, что правая часть есть $\mathcal{C}$-измеримая с.в. Проверим интегральное свойство. Пусть сначала $\eta = \I_B$ для некоторого $B \in \mathcal{C}$. Тогда для любого $A \in \mathcal{C}$
        \begin{align*}
            \E(\eta\xi \cdot \I_A) &= \E(\xi \cdot \I_B \cdot \I_A) = \E(\xi \cdot \I_{A \cap B}) = \\
            &= \E(\E(\xi|\mathcal{C})\I_{A \cap B}) = \E(\E(\xi|\mathcal{C})\I_B \cdot \I_A) = \E(\eta\E(\xi|\mathcal{C})\I_A).
        \end{align*}
        Далее, в силу линейности математического ожидания, интегральное свойство будет выполнено и для простых с.в. вида $\eta = \sum_{k=1}^m c_k \cdot \I_{B_k}$, $c_k \in \mathbb{R}$, $B_k \in \mathcal{C}$. 
        
        Для произвольной $\eta$ возьмем из теоремы 2.9 последовательность простых $\mathcal{C}$-измеримых с.в. $\eta_n$ такую, что $\eta_n \toAS \eta$ и $|\eta_n| \le |\eta|$. Тогда по свойству (8) получаем, что
        $$ \E(\eta_n\xi|\mathcal{C}) \toAS \E(\eta\xi|\mathcal{C}). $$
        Однако
        $$ \E(\eta_n\xi|\mathcal{C}) = \eta_n\E(\xi|\mathcal{C}) \toAS \eta\E(\xi|\mathcal{C}). $$
        Следовательно, $\E(\xi\eta|\mathcal{C}) = \eta\E(\xi|\mathcal{C})$.
        \item В силу выпуклости для любого $x \in \mathbb{R}$ существует $\lambda(x)$ т.,ч. для любого $y \in \mathbb{R}$ выполнено
        $$ \varphi(y) \ge \varphi(x) + \lambda(x)(y - x). $$
        Положим $y = \xi$, $x = \E(\xi|\mathcal{C})$ и возьмем УМО относительно $\mathcal{C}$ от обеих частей неравенства. Пользуясь свойствами (1) и (9), получим:
        \begin{align*}
            \E(\varphi(\xi)|\mathcal{C}) &\ge \E\Big( \varphi(\E(\xi|\mathcal{C})) \mid \mathcal{C} \Big) + \E\Big( \lambda(\E(\xi|\mathcal{C}))(\xi - \E(\xi|\mathcal{C})) \mid \mathcal{C} \Big) = \\
            &= [\text{выносим } \lambda(\E(\xi|\mathcal{C})) \text{ как } \mathcal{C}\text{-измеримую с.в.}] = \\
            &= \varphi(\E(\xi|\mathcal{C})) + \lambda(\E(\xi|\mathcal{C})) \E(\xi - \E(\xi|\mathcal{C}) \mid \mathcal{C}) = \varphi(\E(\xi|\mathcal{C})).
        \end{align*}
    \end{enumerate}
}

Следующее свойство показывает, что при наличии второго момента условное математическое ожидание $\E(\xi|\mathcal{C})$ можно воспринимать как проекцию случайной величины $\xi$ на подпространство $\mathcal{C}$-измеримых случайных величин в $L^2(\Omega, \mathcal{F}, \P)$.

\mkthm{О наилучшем квадратичном прогнозе}{l2_forecast}{
    Пусть $\xi$ --- случайная величина на $(\Omega, \mathcal{F}, \P)$ с конечным вторым моментом, $\mathcal{C}$ --- под-$\sigma$-алгебра в $\mathcal{F}$. Обозначим через $\mathcal{A}_{\mathcal{C}}$ --- множество всех $\mathcal{C}$-измеримых с.в. с конечным вторым моментом. Тогда
    $$ \E(\xi - \E(\xi|\mathcal{C}))^2 = \min_{\eta \in \mathcal{A}_{\mathcal{C}}} \E(\xi - \eta)^2, $$
    причем $\E(\xi|\mathcal{C})$ --- это единственный минимум с точностью до равенства п.н.
}
\mkproof{
    Пусть $\eta \in \mathcal{A}_{\mathcal{C}}$. Тогда
    \begin{align*}
        \E(\xi - \eta)^2 &= \E(\xi - \E(\xi|\mathcal{C}) + \E(\xi|\mathcal{C}) - \eta)^2 = \\
        &= \E(\xi - \E(\xi|\mathcal{C}))^2 + \E(\E(\xi|\mathcal{C}) - \eta)^2 + 2\E\Big( (\xi - \E(\xi|\mathcal{C}))(\E(\xi|\mathcal{C}) - \eta) \Big).
    \end{align*}
    Проверим, что третье слагаемое равно нулю. Действительно, по формуле полной вероятности
    \begin{align*}
        \E\Big( (\xi - \E(\xi|\mathcal{C}))(\E(\xi|\mathcal{C}) - \eta) \Big) &= \E\Big( \E\Big( (\xi - \E(\xi|\mathcal{C}))(\E(\xi|\mathcal{C}) - \eta) \mid \mathcal{C} \Big) \Big) = \\
        \intertext{(свойство (9) теоремы \ref{thm:cme_props})}
        &= \E\Big( (\E(\xi|\mathcal{C}) - \eta) \E(\xi - \E(\xi|\mathcal{C}) \mid \mathcal{C}) \Big) = 0.
    \end{align*}
    В итоге,
    $$ \E(\xi - \eta)^2 = \E(\xi - \E(\xi|\mathcal{C}))^2 + \E(\E(\xi|\mathcal{C}) - \eta)^2 \ge \E(\xi - \E(\xi|\mathcal{C}))^2, $$
    причем равенство достигается тогда и только тогда, когда $\eta = \E(\xi|\mathcal{C})$.
}

Условное математическое ожидание случайной величины относительно сигма-алгебры --- весьма общий объект, позволяющий использовать его в разных ситуациях, например, когда сигма-алгебра порождена некоторым случайным процессом. В следующем параграфе мы рассмотрим наиболее простую ситуацию, когда она порождается другой случайной величиной или вектором.

\subsection*{6.3 Условные распределения}

Пусть $(\Omega, \mathcal{F}, \P)$ --- вероятностное пространство.

\mkdef{Условная вероятность события}{cond_prob_sigma}{
    Пусть $A \in \mathcal{F}$ --- событие. Тогда для под-$\sigma$-алгебры $\mathcal{C} \subset \mathcal{F}$ определим
    $$ \P(A|\mathcal{C}) := \E(\I_A|\mathcal{C}). $$
    Если $\xi$ и $\eta$ --- две случайные величины, то полагаем
    $$ \E(\xi|\eta) := \E(\xi|\mathcal{F}_\eta), $$
    где $\mathcal{F}_\eta$ --- $\sigma$-алгебра, порожденная с.в. $\eta$.
}

\mkrem{
    Напомним, что с.в. $\zeta$ является $\mathcal{F}_\eta$-измеримой тогда и только тогда, когда существует такая борелевская функция $\psi(x)$, что $\zeta = \psi(\eta)$.
}

Возникает естественный вопрос о том, как можно было бы вычислять условные математические ожидания $\E(\xi|\eta)$. Оказывается, ответ будет почти таким же, как и для обычных (безусловных) математических ожиданий. А именно, с помощью интегралов по распределениям, но в этот раз --- условным. Начнем обсуждение условных распределений с определения условного математического ожидания $\E(\xi|\eta=y)$.

\mkdef{УМО случайной величины относительно другой с.в.}{cme_var}{
    Пусть $\xi$ и $\eta$ --- две случайные величины. Тогда величиной $\E(\xi|\eta = y)$ называется такая борелевская функция $\varphi(y)$, что для любого $B \in \mathcal{B}(\mathbb{R})$ выполняется равенство
    $$ \E(\xi \cdot \I\{\eta \in B\}) = \E(\varphi(\eta) \cdot \I\{\eta \in B\}) = \int\limits_B \varphi(y)\P_\eta(dy). $$
}

Отметим, что в отличие от формулы \eqref{eq:cme_discrete} данное определение понимает величину $\E(\xi|\eta = y)$ как борелевскую функцию целиком. При конкретном $y$ ее значение, вообще говоря, может быть произвольным.

\mklem{О существовании}{cme_var_exist}{
    Если $\E|\xi| < +\infty$, то $\E(\xi|\eta = y)$ существует и единственно (с точностью до равенства $\P_\eta$-п.н.).
}
\mkproof{
    Для борелевских множеств $B \in \mathcal{B}(\mathbb{R})$ положим
    $$ Q(B) = \E(\xi \cdot \I\{\eta \in B\}). $$
    Заметим, что это заряд на вероятностном пространстве $(\mathbb{R}, \mathcal{B}(\mathbb{R}), \P_\eta)$, абсолютно непрерывный относительно $\P_\eta$. По теореме Радона--Никодима существует единственная $\P_\eta$-п.н. с.в. $\varphi$ на $(\mathbb{R}, \mathcal{B}(\mathbb{R}), \P_\eta)$ (борелевская функция!) т.,ч. для любого $B \in \mathcal{B}(\mathbb{R})$
    $$ Q(B) = \int\limits_B \varphi(y)\P_\eta(dy). $$
    По определению получаем, что $\varphi(y) = \E(\xi|\eta = y)$.
}

Естественно задаться вопросом как связаны $\E(\xi|\eta)$ и $\E(\xi|\eta = y)$. Несложно понять, что по сути это одно и тоже.

\mkprop{}{cme_prop_equiv}{
    $$ \E(\xi|\eta = y) = \varphi(y) \iff \E(\xi|\eta) = \varphi(\eta). $$
}
\mkproof{
    ($\implies$) Пусть $B \in \mathcal{B}(\mathbb{R})$. Тогда
    $$ \E(\xi \cdot \I\{\eta \in B\}) = \int\limits_B \varphi(y)\P_\eta(dy) = \E(\varphi(\eta) \cdot \I\{\eta \in B\}). $$
    Значит, $\varphi(\eta) = \E(\xi|\eta)$.
    
    ($\impliedby$) То же самое в обратном порядке:
    $$ \E(\xi \cdot \I\{\eta \in B\}) = \E(\varphi(\eta) \cdot \I\{\eta \in B\}) = \int\limits_B \varphi(y)\P_\eta(dy). $$
    Следовательно, $\varphi(y) = \E(\xi|\eta = y)$.
}

Условное математическое ожидание $\E(\xi|\eta = y)$ обладает многими свойствами условного математического ожидания $\E(\xi|\eta)$. В частности, будут выполнены
\begin{itemize}
    \item линейность,
    \item сохранение отношения порядка,
    \item предельный переход под знаком УМО,
    \item неравенство Йенсена.
\end{itemize}
Доказываются они точно так же, как и свойства $\E(\xi|\eta)$. В связи с этим мы опустим их подробное обсуждение.

Перейдем непосредственно к обсуждению условных распределений.

\mkdef{Условное распределение}{cond_dist}{
    Условным распределением случайной величины $\xi$ относительно случайной величины $\eta$ называется такая функция $\P(B, y)$, $B \in \mathcal{B}(\mathbb{R})$, $y \in \mathbb{R}$, что
    \begin{enumerate}
        \item $\P(\cdot, y)$ является распределением вероятностей на $(\mathbb{R}, \mathcal{B}(\mathbb{R}))$ для каждого $y \in \mathbb{R}$;
        \item $\P(B, \cdot)$ является борелевской функцией для каждого $B \in \mathcal{B}(\mathbb{R})$;
        \item Для любых $B \in \mathcal{B}(\mathbb{R})$ и $A \in \mathcal{B}(\mathbb{R})$ выполняется равенство
        $$ \P(\xi \in B, \eta \in A) = \E(\P(B, \eta) \cdot \I\{\eta \in A\}). $$
    \end{enumerate}
}

\mkrem{
    В принципе, для любого $B \in \mathcal{B}(\mathbb{R})$ мы могли бы определить условную вероятность вида
    \begin{equation} \label{eq:cond_prob_var}
        \P(\xi \in B|\eta = y) := \E(\I\{\xi \in B\}|\eta = y).
    \end{equation}
    Несложно проверить, что тогда будут выполнены свойства 2 и 3 из определения \ref{def:cond_dist}, но свойство 1 (грубо говоря, счетная аддитивность по $B$) не будет гарантировано для всех последовательностей множеств одновременно. Из-за этого формально необходимо отдельно доказывать существование условного распределения.
}

\mkthm{Условное распределение существует}{cond_dist_exist}{
    Условное распределение существует.
}

Доказательство теоремы \ref{thm:cond_dist_exist} можно найти в [1]. Всюду далее будем использовать обозначение $\P(\xi \in B|\eta = y)$ для условного распределения из определения \ref{def:cond_dist}. Отметим, что для него будет верна формула \eqref{eq:cond_prob_var}. Как вероятностная мера, условное распределение также может иметь плотность.

\mkdef{Условная плотность}{cond_density}{
    Условной плотностью случайной величины $\xi$ относительно случайной величины $\eta$ называется плотность условного распределения (по некоторой мере $\mu$), т.е. такая функция $f_{\xi|\eta}(x|y)$, $x \in \mathbb{R}$, $y \in \mathbb{R}$, что для любого $B \in \mathcal{B}(\mathbb{R})$ выполняется равенство:
    $$ \P(\xi \in B|\eta = y) = \int\limits_B f_{\xi|\eta}(x|y)\mu(dx). $$
}

С помощью условных плотностей можно находить условные математические ожидания по той же формуле, что и обычные математические ожидания с помощью обычных плотностей.

\mkthm{}{cme_via_density}{
    Если существует условная плотность $f_{\xi|\eta}(x|y)$ случайной величины $\xi$ относительно случайной величины $\eta$, то для любой борелевской функции $g(x) \colon \mathbb{R} \to \mathbb{R}$ выполнено:
    $$ \E(g(\xi)|\eta = y) = \int\limits_{\mathbb{R}} g(x)f_{\xi|\eta}(x|y)\mu(dx). $$
}
\mkproof{
    Пусть сначала $g(x) = \I\{x \in A\}$ для некоторого борелевского множества $A$. Тогда для любого $y$
    \begin{align*}
        \E(g(\xi)|\eta = y) &= \P(\xi \in A|\eta = y) = \int\limits_A f_{\xi|\eta}(x|y)\mu(dx) = \\
        &= \int\limits_{\mathbb{R}} \I\{x \in A\}f_{\xi|\eta}(x|y)\mu(dx) = \int\limits_{\mathbb{R}} g(x)f_{\xi|\eta}(x|y)\mu(dx).
    \end{align*}
    Раз формула выполнена для индикаторов, то в силу линейности математического ожидания формула верна и для линейных комбинаций индикаторов, т.е. для простых функций $g(x)$. Для произвольной функции $g(x)$ справедливость формулы устанавливается предельным переходом от простых функций.
}

Остается, впрочем, вопрос о том, как вычислять саму условную плотность. Оказывается, она всегда есть, если имеется совместная плотность случайных величин.

\mkthm{Достаточное условие существования условной плотности}{cond_density_exist}{
    Пусть случайные величины $\xi$ и $\eta$ имеют совместную плотность $f_{\xi,\eta}(x, y)$ по мере $\mu \times \nu$ на $\mathbb{R}^2$. Тогда функция
    $$ f_{\xi|\eta}(x|y) = \begin{cases} \frac{f_{\xi,\eta}(x,y)}{f_\eta(y)}, & \text{если } f_\eta(y) > 0; \\ p(x), & \text{иначе;} \end{cases} $$
    является условной плотностью $\xi$ относительно $\eta$ по мере $\mu$. Здесь $f_\eta(y)$ --- это плотность случайной величины $\eta$ по мере $\nu$, а $p(x)$ --- произвольная плотность на прямой по мере $\mu$.
}
\mkproof{
    Надо проверить, что функция $\P(B, y) = \int\limits_B f_{\xi|\eta}(x|y)\mu(dx)$ является условным распределением $\xi$ относительно $\eta$. В силу свойств интеграла будет выполнено свойство 1) из определения \ref{def:cond_dist}, а из теоремы Фубини следует, что правая часть будет удовлетворять свойству 2). Остается проверить только свойство 3).
    
    Для любых борелевских множеств $A, B$ выполнено
    \begin{align*}
        \P(\xi \in B, \eta \in A) &= \int\limits_{B \times A} f_{\xi,\eta}(x, y)\mu(dx)\nu(dy) = \\
        &= \int\limits_A \left( \int\limits_B \frac{f_{\xi,\eta}(x, y)}{f_\eta(y)}\mu(dx) \right) f_\eta(y)\nu(dy) = \int\limits_A \P(B, y)f_\eta(y)\nu(dy) = \\
        &= \int\limits_{\mathbb{R}} \P(B, y) \cdot \I\{y \in A\}f_\eta(y)\nu(dy) = \E(\P(B, \eta) \cdot \I\{\eta \in A\}).
    \end{align*}
    Последнее равенство означает, что функция $\frac{f_{\xi,\eta}(x,y)}{f_\eta(y)}$ является условной плотностью.
}

Теоремы \ref{thm:cond_dist_exist} и \ref{thm:cme_via_density} позволяют сформулировать следующую схему вычисления условного математического ожидания. Допустим мы хотим вычислить $\E(g(\xi)|\eta)$. Тогда
\begin{enumerate}
    \item Находим $f_{\xi,\eta}(x, y)$ --- совместную плотность $\xi$ и $\eta$.
    \item С ее помощью находим $f_{\xi|\eta}(x|y)$ --- условную плотность $\xi$ относительно $\eta$.
    \item Вычисляем
    $$ h(y) = \E(g(\xi)|\eta = y) = \int\limits_{\mathbb{R}} g(x)f_{\xi|\eta}(x|y)\mu(dx). $$
    \item Подставляем $\E(g(\xi)|\eta) = h(\eta)$.
\end{enumerate}

Проиллюстрируем приведенную схему примером.

\vspace{1em}
\noindent\textbf{Упражнение 6.1.} Случайные величины $X$ и $Y$ независимые с распределением $\text{Exp}(1)$. Найдите
$$ \E(X^2|X + Y) = ? $$

\noindent\textbf{Решение.} Сначала найдем совместную плотность $(X, X + Y)$. Для $B \in \mathcal{B}(\mathbb{R}^2)$ рассмотрим
\begin{align*}
    \P((X, X + Y) \in B) &= \int\limits_{(x,y): (x,x+y) \in B} f_{X,Y}(x, y) \, dxdy = \\
    &= \int\limits_{(x,y): (x,x+y) \in B} e^{-x} e^{-y} \cdot \I\{x, y > 0\} \, dxdy = \\
    \intertext{(сделаем замену $\alpha = x, \beta = x + y$, якобиан замены $= 1$)}
    &= \int\limits_{(\alpha,\beta) \in B} e^{-\beta} \cdot \I\{\beta > \alpha > 0\} \, d\alpha d\beta.
\end{align*}
Стало быть, функция $p_{X,X+Y}(\alpha, \beta) = e^{-\beta} \cdot \I\{\beta > \alpha > 0\}$ и есть совместная плотность $(X, X + Y)$. Теперь найдем плотность $X + Y$. Для этого совместную плотность $p_{X,X+Y}(\alpha, \beta)$ проинтегрируем по первой переменной.
\begin{align*}
    p_{X+Y}(\beta) &= \int\limits_{\mathbb{R}} p(\alpha, \beta) \, d\alpha = \int\limits_{\mathbb{R}} e^{-\beta} \I\{\beta > \alpha > 0\} \, d\alpha = \\
    &= e^{-\beta} \I\{\beta > 0\} \int\limits_0^\beta \, d\alpha = \beta e^{-\beta} \I\{\beta > 0\}.
\end{align*}
Отсюда условная плотность равна
$$ f_{X|X+Y}(\alpha|\beta) = \frac{p_{X,X+Y}(\alpha, \beta)}{p_{X+Y}(\beta)} = \frac{1}{\beta} \I\{\beta > \alpha > 0\}. $$
Остается найти УМО
$$ \E(X^2|X + Y = \beta) = \int\limits_{\mathbb{R}} \alpha^2 \cdot \frac{1}{\beta} \I\{\beta > \alpha > 0\} \, d\alpha = \frac{\beta^2}{3}. $$
В итоге, $\E(X^2|X + Y) = (X + Y)^2 / 3$. \hfill $\square$
\vspace{1em}

\mkrem{
    Сделаем еще несколько замечаний.
    \begin{itemize}
        \item Если $\xi$ и $\eta$ независимы и имеют плотности, то условная плотность $f_{\xi|\eta}(x|y)$ равна просто плотности случайной величины $\xi$.
        \item Вычисление $\E(\xi|\eta)$ в случае векторного условия $\eta$ происходит по тем же формулам.
        \item Если условной плотности нет, то формально условное математическое ожидание вычисляется как интеграл по условному распределению
        $$ \E(g(\xi)|\eta = y) = \int\limits_{\mathbb{R}} g(x)\mathcal{Q}_y(dx), $$
        где $\mathcal{Q}_y(B) = \P(\xi \in B|\eta = y)$.
        \item Особо важно уметь вычислять условные математические ожидания вида $\E(g(\xi, \eta)|\eta)$ для независимых случайных величин $\xi$ и $\eta$.
    \end{itemize}
}

В последнем случае есть очень полезная формула, позволяющая быстро вычислять условные математические ожидания.

\mklem{}{cme_indep}{
    Если $\xi$ и $\eta$ --- независимые случайные величины, то
    $$ \E(g(\xi, \eta)|\eta = y) = \E g(\xi, y). $$
}
\mkproof{
    По теореме Фубини правая часть есть борелевская функция от $y$. Проверим интегральное свойство: для любого $B \in \mathcal{B}(\mathbb{R})$
    $$ \E(g(\xi, \eta) \cdot \I\{\eta \in B\}) = \int\limits_{\mathbb{R}^2} g(x, y) \cdot \I\{y \in B\} \P_{\xi,\eta}(dx, dy) = $$
    (пользуемся независимостью и снова теоремой Фубини)
    $$ = \int\limits_B \left( \int\limits_{\mathbb{R}} g(x, y)\P_\xi(dx) \right) \P_\eta(dy) = \int\limits_B \E g(\xi, y)\P_\eta(dy). $$
}

